\mychapter{Plano de Continuidade}{cap:plano_continuidade}\lhead{PLANO DE CONTINUIDADE}

A continuidade parte de uma avaliação temporal da Bahia já concluída e de uma avaliação M5 em andamento. O objetivo é consolidar evidência comparável para a seleção de políticas e, então, avaliar o PSE supervisionado. A expansão nacional foi cancelada na execução atual e não constitui requisito para concluir as etapas seguintes.

\section{Consolidação dos experimentos}

A Bahia dispõe de manifesto, 4.350 ajustes completos e métricas de validação e teste. A análise deverá examinar a variação entre origens, os casos sem demanda e o compromisso entre custo e atendimento. As análises complementam a mesma matriz experimental, com cobertura e proveniência verificadas.

\subsection{M5: avaliação pública com protocolo temporal}
\label{sec:estudos_caso_adicionais}

A execução M5 utiliza a coorte histórica de 2.408 identificadores, sujeita à elegibilidade no prefixo de cada origem. Isso preserva uma comparação condicional à coorte, sem tornar prospectiva sua seleção inicial. O recorte local é inspirado em \citet{Zabraoui2025GaDQN}; não reproduz uma amostra publicada com identificadores e filtros numéricos integralmente especificados.

A comparação com a Bahia deverá usar métricas e regras de escolha comuns, mantendo separadas as escalas de custo e as propriedades da demanda. A avaliação M5 ainda não sustenta conclusões neste documento.

\section{Validação da seleção}

O próximo desenho deverá ampliar os períodos e reduzir a sobreposição dos testes. Características, normalização, parâmetros e critérios de inclusão devem depender do passado disponível. A seleção usará validação, e o teste ficará reservado à avaliação. O comparador principal será a melhor política global viável escolhida sob a mesma informação.

O PSE receberá características anteriores à decisão e rótulos produzidos apenas nos dados de ajuste. Sua avaliação incluirá custo, atendimento, déficit de serviço e arrependimento relativo à referência. Acurácia e F1 serão diagnósticos auxiliares: rótulos próximos em desempenho não têm o mesmo significado econômico que erros com grande perda de serviço. Modelos supervisionados serão comparados com política global, regras por perfil e modelos simples, respeitando tempo e agrupamento por loja.

\subsection{Calibração agrupada por perfil}
\label{sec:robustez_pooling}

Calibrar parâmetros em grupos é uma hipótese distinta de escolher políticas por perfil. Uma eventual comparação entre ajuste individual e agrupado deverá manter cortes, custos e orçamento comuns. Benefícios e inversões entre períodos serão apresentados conjuntamente, sem interpretar um recorte favorável como robustez geral.

\section{Sensibilidade e eficiência}

As análises deverão variar custos, prazo, estoque inicial, tratamento terminal e exigência de serviço. Média por série, atendimento ponderado e restrições individuais serão comparados como problemas de decisão distintos. A escolha de um peso para custo e serviço deverá anteceder o teste, evitando ajustar o objetivo para favorecer um resultado observado.

A revisão do portfólio combinará diversidade das regras, custo computacional e evidência nas duas bases. Métodos semelhantes não serão considerados equivalentes apenas por correlação de custo. A eventual retirada de candidatos será definida nos dados de desenvolvimento e documentada antes da avaliação final. GA--DQN continuará sendo comparado com GA e DQN, com o esforço total da combinação informado.

CPU e GPU serão escolhidas por medições do tempo total e do uso de recursos. A execução M5 atual usa oito processos CPU para coexistir com o uso normal da máquina. Essa configuração operacional não altera sementes, cortes ou candidatos, nem é um resultado de qualidade das políticas.

\section{Consolidação do texto e critérios de conclusão}

O artigo para \textit{Supply Chain Analytics} e a dissertação compartilham resultados conferidos, mas têm funções diferentes: o artigo concentra a avaliação incremental da seleção; a dissertação explica fundamentos, percurso e limites. O material SBPO permanece arquivado. Autoria, declarações, disponibilidade de dados e adequação editorial serão consolidadas antes de submissão ou depósito.

\begin{table}[htbp]
\centering\small
\caption{Etapas de continuidade e critérios de conclusão, sem datas ainda não acordadas.}
\label{tab:cronograma}
\begin{tabular}{@{}p{4.1cm}p{8.9cm}@{}}\toprule
Etapa & Critério de conclusão\\\midrule
Consolidação Bahia & Resultados temporais, falhas de serviço e limites documentados\\
Conclusão M5 & Cobertura planejada, matriz completa e comparação por origem\\
Seleção e sensibilidade & Critérios de decisão definidos e teste preservado\\
Avaliação do PSE & Comparação fora da amostra com referências competitivas\\
Redação e defesa & Afirmações, contribuições e pendências coerentes com a evidência\\
\bottomrule\end{tabular}
\end{table}
